%% Đề 43 -- de-tuyen-sinh-lop-10-thpt-mon-toan-nam-2015-2016-so-gddt-lang-son.pdf
\deheader{ĐỀ SỐ 43}%
  {Kỳ thi Tuyển sinh vào lớp 10 THPT — Năm học 2015--2016}%
  {Sở GD\&ĐT Lạng Sơn | Chỉ phần Đại số}
\addcontentsline{toc}{section}{Đề 43 — Tuyển sinh vào lớp 10 THPT — Năm học 2015--2016}

\begin{de_cau}{Câu 1}{2,0 điểm}
\begin{enumerate}[label=\alph*)]
  \item Tính: $A = \dfrac{1}{\sqrt{5}-2}+\sqrt{20}$.
  \item Rút gọn: $B = \left(\dfrac{1}{\sqrt{x}+1}-\dfrac{1}{\sqrt{x}}\right):\dfrac{1}{x+\sqrt{x}}$ với $x>0$.
\end{enumerate}
\end{de_cau}

\begin{loi_giai_cau}
\textbf{a)} $\dfrac{1}{\sqrt{5}-2}=\dfrac{\sqrt{5}+2}{(\sqrt{5})^2-4}=\sqrt{5}+2$. $\sqrt{20}=2\sqrt{5}$.

$A=\sqrt{5}+2+2\sqrt{5}=\boxed{3\sqrt{5}+2}$.

\textbf{b)} $\dfrac{1}{\sqrt{x}+1}-\dfrac{1}{\sqrt{x}}=\dfrac{\sqrt{x}-(\sqrt{x}+1)}{\sqrt{x}(\sqrt{x}+1)}=\dfrac{-1}{\sqrt{x}(\sqrt{x}+1)}$.

$x+\sqrt{x}=\sqrt{x}(\sqrt{x}+1)$, nên $\dfrac{1}{x+\sqrt{x}}=\dfrac{1}{\sqrt{x}(\sqrt{x}+1)}$.

$B=\dfrac{-1}{\sqrt{x}(\sqrt{x}+1)}\cdot\sqrt{x}(\sqrt{x}+1)=\boxed{-1}$.
\end{loi_giai_cau}

\begin{de_cau}{Câu 2}{2,0 điểm}
\begin{enumerate}[label=\alph*)]
  \item Giải phương trình: $x^2-7x+10=0$.
  \item Giải hệ phương trình:
  \[\begin{cases} 3x-2y=1 \\ x+3y=11 \end{cases}.\]
\end{enumerate}
\end{de_cau}

\begin{loi_giai_cau}
\textbf{a)} $(x-2)(x-5)=0\Rightarrow\boxed{x=2}$ hoặc $\boxed{x=5}$.

\textbf{b)} Nhân pt(1) với 3: $9x-6y=3$; nhân pt(2) với 2: $2x+6y=22$. Cộng: $11x=25\Rightarrow x=\dfrac{25}{11}$... Hoặc: $x=11-3y$ thế vào (1): $3(11-3y)-2y=1\Rightarrow33-11y=1\Rightarrow y=\dfrac{32}{11}$.

Vậy $(x;y)=\boxed{\left(\dfrac{25}{11};\dfrac{32}{11}\right)}$.
\end{loi_giai_cau}

\begin{de_cau}{Câu 3}{2,0 điểm}
Cho phương trình $x^2-2mx+m^2-4m=0$ (tham số $m$).
\begin{enumerate}[label=\alph*)]
  \item Giải phương trình khi $m=5$.
  \item Tìm $m$ để phương trình có hai nghiệm dương phân biệt.
\end{enumerate}
\end{de_cau}

\begin{loi_giai_cau}
\textbf{a)} $m=5$: $x^2-10x+5=0\Rightarrow x=5\pm2\sqrt{5}$.

\textbf{b)} $\Delta'=m^2-(m^2-4m)=4m>0\Rightarrow m>0$.

Vi-ét: $x_1+x_2=2m>0$ (\checkmark  nếu $m>0$) và $x_1x_2=m^2-4m>0\Rightarrow m(m-4)>0\Rightarrow m>4$ (kết hợp $m>0$).

Vậy $\boxed{m>4}$.
\end{loi_giai_cau}

\begin{de_cau}{Câu 5}{0,5 điểm}
Cho $a, b, c > 0$ thỏa mãn $a + b + c = 3$. Chứng minh:
\[\sqrt{a} + \sqrt{b} + \sqrt{c} \le \sqrt{3(a+b+c)} = 3.\]
Tìm điều kiện đẳng thức.
\end{de_cau}

\begin{loi_giai_cau}
Theo Cauchy-Schwarz: $(\sqrt{a}+\sqrt{b}+\sqrt{c})^2\le(1+1+1)(a+b+c)=3\cdot3=9$.

Nên $\sqrt{a}+\sqrt{b}+\sqrt{c}\le3$. Dấu $=$ khi $a=b=c=1$. $\square$
\end{loi_giai_cau}
